% Part: normal-modal-logic
% Chapter: axioms-systems
% Section: duals

\documentclass[../../../include/open-logic-section]{subfiles}

\begin{document}

\olfileid[pt]{nml}{prf}{dua}

\olsection{\usetoken{P}{formula} duais}

\begin{defn}\ollabel{def:duals}
  Cada uma das !!{formula}s \Ax{T}, \Ax{B}, \Ax{4} e
  \Ax{5} tem uma \emph{dual}, denotada por um losango subscrito, como
  segue:
  \begin{align}
    \tag{\Ax{T_\Diamond}} p & \lif \Diamond p\\
    \tag{\Ax{B_\Diamond}} \Diamond\Box p & \lif p\\
    \tag{\Ax{4_\Diamond}} \Diamond\Diamond p & \lif \Diamond p\\
    \tag{\Ax{5_\Diamond}} \Diamond\Box p & \lif \Box p
    \end{align}
\end{defn}

Cada uma das !!{formula}s duais acima é obtida da !!{formula}
correspondente substituindo $\lnot p$ por $p$, tomando a contrapositiva,
substituindo $\lnot\Box\lnot$ por $\Diamond$ e substituindo
$\lnot\Diamond\lnot$ por~$\Box$. \Ax{D}, O.s.\@,
$\Box!A \lif \Diamond!A$, é sua própria dual nesse sentido.

\begin{prop}\ollabel{prop:dualsys}
  Para cada !!{formula}~$!A$ em \olref[nml][prf][dua]{def:duals}:
  $\Log{K}!A = \Log{K}!A_{\Diamond}$.
\end{prop}
\begin{proof}
  Exercício.
\end{proof}
\begin{prob}
  Demonstre \olref[nml][prf][dua]{prop:dualsys}.
\end{prob}

\end{document}
